git\documentclass[a4paper,12pt]{article}
\usepackage[T2A]{fontenc}
\usepackage[utf8]{inputenc}
\usepackage[russian]{babel}
\usepackage{amsmath,amssymb,mathtools}
\usepackage{PTSans} % кириллический академический sans-serif, pdfLaTeX
\renewcommand{\familydefault}{\sfdefault}
\usepackage{icomma}
\usepackage[margin=22mm,headheight=26pt,headsep=6mm,footskip=12mm]{geometry}
\usepackage{microtype}
\usepackage{xcolor}
\usepackage{tikz}
\usetikzlibrary{arrows.meta,positioning,shapes.geometric,calc}
\usepackage{booktabs,tabularx}
\usepackage{enumitem}
\usepackage{parskip}
\usepackage{fancyhdr}
\usepackage{setspace}
\usepackage{array}

\definecolor{accent}{HTML}{008080}
\definecolor{darkaccent}{HTML}{004D4D}
\definecolor{textgray}{HTML}{333333}
\definecolor{softgray}{HTML}{6F7777}
\color{textgray}
\onehalfspacing
\setlength{\parindent}{0pt}
\setlength{\parskip}{0.45em}
\setlist[itemize]{leftmargin=*,topsep=3pt,itemsep=3pt,parsep=0pt}
\setlist[enumerate]{leftmargin=*,topsep=4pt,itemsep=6pt,parsep=0pt}
\renewcommand{\arraystretch}{1.25}

% Дата занятия, фиксированная для архивной версии документа.
\year=2026 \month=10 \day=25

\pagestyle{fancy}
\fancyhf{}
\fancyhead[L]{\small\color{softgray}Математика \textbullet{} 7 класс}
\fancyhead[R]{\small\color{softgray}25.10.2026}
\fancyfoot[C]{\small\color{softgray}\thepage}
\renewcommand{\headrulewidth}{0.3pt}
\renewcommand{\headrule}{\hbox to\headwidth{\color{accent!55}\leaders\hrule height \headrulewidth\hfill}}
\renewcommand{\footrulewidth}{0pt}

\newcommand{\sectionline}{\par\vspace{-0.65em}\noindent{\color{accent!65}\rule{\linewidth}{0.45pt}}\par\vspace{0.25em}}
\newcommand{\knowledge}[1]{\noindent$\square$\enspace #1\par}
\newcommand{\data}[1]{\left(#1\right)}

\begin{document}
\thispagestyle{empty}
\vspace*{37mm}
\begin{center}
{\color{darkaccent}\large УЧЕБНОЕ ПОСОБИЕ \textbullet{} 7 КЛАСС}\par
\vspace{18mm}
{\color{darkaccent}\fontsize{29}{35}\selectfont\bfseries Чек-лист}\par
\vspace{7mm}
{\fontsize{22}{28}\selectfont\bfseries Среднее арифметическое,\\[3mm] медиана и размах}\par
\vspace{12mm}
{\large Числовой ряд и три способа его описания}\par
\vspace{22mm}
{\color{accent}\rule{44mm}{0.7pt}}\par
\vspace{8mm}
{\large \today}\par
\vspace{18mm}
{\large Лёвин Артём Александрович}\par
{\large\color{darkaccent}эксклюзивно для Кирилла Зиновьева}\par
\end{center}
\vfill
\begin{center}
\begin{tikzpicture}[x=1cm,y=1cm]
\draw[accent!55,line width=1.15pt] (0,0)
  .. controls (2.5,0.95) and (4.5,-0.8) .. (7,0)
  .. controls (9.5,0.8) and (11.5,-0.95) .. (14,0);
\end{tikzpicture}
\end{center}
\clearpage

\section*{1. Подготовка числового ряда}
\sectionline
\knowledge{\textbf{Числовой ряд} --- набор наблюдаемых значений; одинаковые числа могут повторяться.}
\knowledge{Для нахождения медианы расположите числа \textbf{по неубыванию}: от меньшего к большему.}

\textbf{Пример.} Исходные данные: $5,10,3,2,15,1$. После упорядочивания:
\[
1\le 2\le 3\le 5\le 10\le 15.
\]
Порядок чисел не влияет на среднее арифметическое и размах, но для медианы он важен.

\section*{2. Среднее арифметическое}
\sectionline
\knowledge{Сложите все $n$ чисел и разделите сумму на их количество:}
\[
\boxed{\displaystyle \overline{x}=\frac{x_1+x_2+\ldots+x_n}{n}=\frac{S}{n}},
\qquad n\ge 1.
\]
\knowledge{Если известны среднее и количество чисел, можно восстановить \textbf{сумму}:}
\[
\boxed{S=n\overline{x}}.
\]

\textbf{Пример 1.} Для ряда $1,2,3,5,10,15$:
\[
\overline{x}=\frac{1+2+3+5+10+15}{6}=\frac{36}{6}=6.
\]

\textbf{Пример 2.} Среднее четырёх чисел $8,10,12,x$ равно $11$.
\[
\frac{8+10+12+x}{4}=11
\quad\Longrightarrow\quad 30+x=44
\quad\Longrightarrow\quad \boxed{x=14}.
\]

\knowledge{Для \textbf{разных неизвестных чисел} используйте индексы $x_1,x_2,\ldots$. По одному известному среднему обычно можно найти их сумму, но нельзя восстановить каждое число.}
\knowledge{Сокращать можно \textbf{множители} дроби. Слагаемые сокращать нельзя:}
\[
\frac{15x}{5}=3x, \qquad \frac{15+x}{5}=3+\frac{x}{5}.
\]

\clearpage
\section*{3. Обратные и составные задачи}
\sectionline
\knowledge{Если среднее $5$ чисел равно $13$, то сумма этих чисел:}
\[
S=5\cdot13=\boxed{65}.
\]
\knowledge{Если известно среднее $n$ наблюдений, сначала найдите их сумму, затем добавьте новые данные и разделите на новое количество наблюдений.}

\textbf{Пример.} За $6$ дней изготовляли в среднем по $20$ деталей в день. За седьмой день изготовили $27$ деталей. Найдём среднюю выработку за все $7$ дней.
\[
S_6=6\cdot20=120,\qquad
\overline{x}_7=\frac{120+27}{7}=\frac{147}{7}=\boxed{21}.
\]
\knowledge{Если числа имеют вид $x,2x,3x,4x,5x$, а их среднее равно $100$, то}
\[
\frac{x+2x+3x+4x+5x}{5}=100
\ \Longrightarrow\ 3x=100
\ \Longrightarrow\ \boxed{x=\frac{100}{3}}.
\]
В таком случае требуется подставить найденное $x$ \textbf{во все пять выражений}, если задача просит найти каждое число.

\section*{4. Медиана}
\sectionline
\knowledge{\textbf{Медиана} $M$ --- центральное значение упорядоченного ряда.}
\knowledge{При \textbf{нечётном} количестве чисел возьмите единственное центральное число.}
\[
3,\ 5,\ \underbrace{7}_{M},\ 8,\ 12
\qquad\Longrightarrow\qquad \boxed{M=7}.
\]
\knowledge{При \textbf{чётном} количестве чисел найдите среднее двух центральных чисел.}
\[
1,\ 2,\ \underbrace{3,\ 4}_{\text{середина}},\ 5,\ 100
\quad\Longrightarrow\quad M=\frac{3+4}{2}=\boxed{3{,}5}.
\]
\knowledge{При повторяющихся значениях правило остаётся тем же. Центральное значение может встречаться несколько раз.}

\clearpage
\section*{Алгоритм: как найти медиану}
\sectionline
\vspace{14mm}
\begin{center}
\begin{tikzpicture}[
  >=Latex,
  node distance=10mm,
  terminal/.style={ellipse,draw=darkaccent,line width=.9pt,minimum width=40mm,minimum height=10mm,align=center,fill=accent!4},
  process/.style={rectangle,rounded corners=2mm,draw=accent,line width=.9pt,minimum width=59mm,minimum height=14mm,align=center,inner sep=3.5mm,text width=51mm},
  test/.style={diamond,aspect=2.1,draw=darkaccent,line width=.9pt,align=center,inner sep=1.5mm,text width=24mm},
  arr/.style={->,line width=.9pt,draw=darkaccent}
]
\node[terminal] (start) {Начало};
\node[process,below=12mm of start] (sort) {Расположите числа\\по неубыванию};
\node[test,below=14mm of sort] (parity) {$n$ нечётно?};
\node[process,below left=21mm and 10mm of parity] (odd) {Возьмите одно\\центральное число};
\node[process,below right=21mm and 10mm of parity] (even) {Найдите среднее\\двух центральных чисел};
\coordinate (join) at ($(odd.south)!0.5!(even.south)+(0,-14mm)$);
\node[process,below=15mm of join] (result) {Запишите медиану $M$};
\node[terminal,below=11mm of result] (end) {Конец};
\draw[arr] (start) -- (sort);
\draw[arr] (sort) -- (parity);
\draw[arr] (parity.west) -| node[pos=.25,above]{\small да} (odd.north);
\draw[arr] (parity.east) -| node[pos=.25,above]{\small нет} (even.north);
\draw[arr] (odd.south) |- (result.west);
\draw[arr] (even.south) |- (result.east);
\draw[arr] (result) -- (end);
\end{tikzpicture}
\end{center}
\clearpage

\section*{5. Размах}
\sectionline
\knowledge{\textbf{Размах} $R$ --- разность наибольшего и наименьшего значений:}
\[
\boxed{R=x_{\max}-x_{\min}},\qquad \boxed{R\ge 0}.
\]
\textbf{Пример.} Для чисел $-11,-8,-5,-2$ размах равен
\[
R=(-2)-(-11)=\boxed{9}.
\]
\knowledge{Если $R=0$, все значения равны. Отрицательным размах быть не может.}
\knowledge{Если величины заданы в разных единицах измерения, сначала переведите их в одни единицы. Например, $2$~м $=200$~см.}

\section*{6. Почему полезны все три характеристики}
\sectionline
Рассмотрим добавление одного очень большого значения.
\begin{center}
\begin{tabularx}{\linewidth}{@{}>{\raggedright\arraybackslash}Xcc@{}}
\toprule
\textbf{Ряд} & \textbf{Среднее} & \textbf{Медиана}\\
\midrule
$2,3,4,5$ & $3{,}5$ & $3{,}5$\\
$2,3,4,5,100$ & $22{,}8$ & $4$\\
\bottomrule
\end{tabularx}
\end{center}
\knowledge{Среднее арифметическое учитывает величину \emph{каждого} числа; оно чувствительно к очень большим и очень маленьким значениям.}
\knowledge{Медиана обычно устойчивее к отдельным крайним значениям, поскольку зависит прежде всего от центральных позиций.}
\knowledge{Размах показывает расстояние между крайними значениями и дополняет среднее с медианой.}

\textbf{Важно.} Одинаковая медиана ещё не означает, что два ряда одинаковы. По одному среднему, медиане или размаху обычно невозможно восстановить все значения ряда.

\textbf{Самопроверка.} Учтите все элементы ряда, включая повторы; упорядочите числа перед поиском медианы. Для чётного количества значений возьмите два центральных. Проверьте знак размаха и единицы измерения.

\clearpage
\section*{Тренировочный блок \textbullet{} Путь А}
\sectionline
\textbf{10 упражнений.} Выполните задания по порядку: от базовых вычислений к сравнению наборов. В заданиях с несколькими вопросами приведите все требуемые результаты.

\begin{enumerate}[label=\textbf{\arabic*.},itemsep=9pt,topsep=5pt]
\item Расположите по неубыванию числа: $8,\ 2,\ 5,\ 2,\ 11$.
\item Найдите среднее арифметическое чисел $4,\ 6,\ 7,\ 9,\ 14$.
\item Найдите медиану ряда $9,\ 4,\ 10,\ 1,\ 4$.
\item Найдите медиану ряда $8,\ 2,\ 5,\ 11,\ 7,\ 3$.
\item Найдите размах чисел $-11,\ -2,\ -8,\ -5$.
\item Среднее арифметическое шести чисел равно $12$. Найдите их сумму.
\item Среднее арифметическое чисел $7,\ 11,\ 13,\ x$ равно $12$. Найдите $x$.
\item Среднее арифметическое пяти чисел $x,\ 2x,\ 3x,\ 4x,\ 5x$ равно $18$. Найдите $x$ и все пять чисел.
\item За пять дней мастер изготовлял в среднем по $14$ деталей в день, а в шестой день изготовил $20$ деталей. Найдите среднюю дневную выработку за все шесть дней.
\item Даны ряды $A=(2,3,5,7,22)$ и $B=(2,3,5,7,7)$. Найдите для каждого ряда среднее арифметическое, медиану и размах. Укажите, какая из трёх характеристик осталась одинаковой.
\end{enumerate}

\clearpage
\section*{Ответы к тренировочному блоку}
\sectionline
\begin{center}
\begin{tabularx}{\linewidth}{@{}>{\bfseries}p{13mm}X@{}}
\toprule
№ & Ответ\\
\midrule
1 & $2,\ 2,\ 5,\ 8,\ 11$.\\
2 & $8$.\\
3 & $4$.\\
4 & $6$.\\
5 & $9$.\\
6 & $72$.\\
7 & $x=17$.\\
8 & $x=6$; числа $6,\ 12,\ 18,\ 24,\ 30$.\\
9 & $15$ деталей в день.\\
10 & $A$: среднее $7{,}8$, медиана $5$, размах $20$.\\
 & $B$: среднее $4{,}8$, медиана $5$, размах $5$.\\
 & Одинакова медиана.\\
\bottomrule
\end{tabularx}
\end{center}

\vspace{8mm}
\textcolor{darkaccent}{\textbf{Контрольная мысль.}} Среднее характеризует общий уровень чисел, медиана --- положение центра упорядоченного ряда, размах --- расстояние между крайними значениями.

\end{document}
